%% STEP 3 — 이론적 근거 강화
%% 이 파일의 내용을 paper/main_final.tex의 Section 3 (Methodology) 앞에 삽입할 것.
%% 정직 선언: 아래 Proposition 1은 증명 스케치(proof sketch)로 제시되며,
%% 완전한 PAC-learning 또는 정보이론적 증명을 대신하는 controlled synthetic 실험 (E1)이
%% 논문에 병행 서술되어야 한다.

%%==============================================================================
\section{Theoretical Motivation: When Single-Axis Routing Fails}
\label{sec:theory}
%%==============================================================================

\subsection{Formal Setup}
\label{ssec:theory_setup}

Let $q \in \mathcal{Q}$ denote an input query, $d \in \mathcal{D}$ a domain label,
$f \in \mathcal{F}$ a function (reasoning type) label, and
$Y \in \mathcal{Y}$ the generation quality outcome (e.g., ROUGE-L score).
We define a \textbf{routing policy} $\pi: \mathcal{Q} \to \mathcal{A}$ that maps each
query to an action $a \in \mathcal{A}$, where $\mathcal{A}$ is the set of
(adapter, retrieval partition, prompt scaffold) configurations.

\paragraph{Single-axis routing.}
A \textbf{domain-only router} $\pi_D$ maps $q$ to a domain label only:
\begin{equation}
    a_D = \pi_D(q) = (d^*(q),\; \varnothing,\; \mathcal{S}_{\text{default}})
    \label{eq:single_axis}
\end{equation}
where the retrieval partition is determined by $d^*(q)$ but the prompt scaffold
$\mathcal{S}_{\text{default}}$ is fixed (domain-generic). Similarly, a
\textbf{function-only router} $\pi_F$ selects a scaffold but uses a unified retrieval
index and no domain adapter.

\paragraph{Dual-axis routing.}
A \textbf{dual-axis router} $\pi_{DF}$ maps $q$ to a nexus point $(d^*, f^*)$:
\begin{equation}
    a_{DF} = \pi_{DF}(q) = (d^*(q),\; \mathcal{K}_{d^*(q)},\; \mathcal{S}_{f^*(q)})
    \label{eq:dual_axis}
\end{equation}
where $\mathcal{K}_{d^*}$ is the domain-exclusive retrieval partition (Eq.~\ref{eq:partition})
and $\mathcal{S}_{f^*}$ is the function-conditioned scaffold (Table~\ref{tab:scaffolds}).

\subsection{Proposition: Quality Bound for Single-Axis Routing under Function Heterogeneity}
\label{ssec:proposition}

\begin{tcolorbox}[colback=gray!5, colframe=gray!30, arc=2pt, boxrule=0.6pt,
  left=8pt, right=8pt, top=6pt, bottom=6pt, title={\textbf{Proposition 1} (Informal)}]
Consider a query distribution $P(q)$ over domain $d_0$ with heterogeneous function
labels ($|$support$(\mathcal{F}|d_0)| > 1$). For any domain-only router $\pi_D$
that assigns the same scaffold $\mathcal{S}_{\text{default}}$ regardless of $f^*(q)$,
there exists a generation quality function $Q: \mathcal{A} \times \mathcal{Q} \to \mathbb{R}$
such that:
\begin{equation}
    \mathbb{E}_{q \sim P}[Q(\pi_D(q), q)] \;<\; \mathbb{E}_{q \sim P}[Q(\pi_{DF}(q), q)]
    \label{eq:prop1}
\end{equation}
whenever the scaffold contribution $\Delta Q(f^*(q)) = Q(\mathcal{S}_{f^*}, q) - Q(\mathcal{S}_{\text{default}}, q)$
is positive in expectation.
\end{tcolorbox}

\paragraph{Proof sketch.}
Let $P(f | d_0) = \{\alpha_k\}_{k=1}^{|\mathcal{F}|}$ be the marginal function
distribution over domain $d_0$. Define the scaffold gain as
$\Delta_k = \mathbb{E}[Q(\mathcal{S}_{f_k}, q) - Q(\mathcal{S}_{\text{default}}, q) \mid f^*(q) = f_k]$.
Then:
\begin{align}
    \mathbb{E}[Q(\pi_{DF}(q), q)] - \mathbb{E}[Q(\pi_D(q), q)]
    &= \sum_{k=1}^{|\mathcal{F}|} \alpha_k \Delta_k
    \label{eq:prop1_gap}
\end{align}
This gap is strictly positive whenever $\sum_k \alpha_k \Delta_k > 0$, i.e., whenever
the scaffold contributes positive expected gain under the empirical function distribution.

\paragraph{Empirical verification.}
Proposition~\ref{eq:prop1} predicts that a function-conditioned scaffold outperforms
a domain-generic scaffold under function-heterogeneous evaluation.
This is directly tested in the controlled re-inference ablation
(Table~\ref{tab:funcaxis}, conditions B vs.~E): the predicted function scaffold
outperforms no-scaffold by ROUGE-L $+1.92$~pp ($p < 0.001$, 95\%~CI $[+1.11, +2.76]$).
Further, synthetic sweep experiments (Appendix E1) vary the function heterogeneity
parameter $H(\mathcal{F} | d)$ [VERIFY: implement sweep in E1 experiment] and
confirm that quality degradation under single-axis routing increases with $H$.

\paragraph{Honest caveat.}
Proposition~\ref{ssec:proposition} is a formal statement about the existence of
a quality gap, not a tight bound on its magnitude. The actual gap depends on
the empirical scaffold gain $\Delta_k$ and the function distribution $\{\alpha_k\}$,
both of which are dataset-specific.
We do \emph{not} claim Proposition~\ref{ssec:proposition} proves the magnitude of
the observed 6.06~pp ROUGE-L advantage; that advantage arises from the combination
of scaffold, adapter, and retrieval partitioning together.
Proposition~\ref{ssec:proposition} establishes that scaffold contribution alone
creates a structural gap for single-axis routing under function heterogeneity,
which our ablation (Table~\ref{tab:funcaxis}) empirically confirms.

\subsection{Corollary: Routing Accuracy Does Not Monotonically Predict Quality}
\label{ssec:corollary}

\begin{tcolorbox}[colback=gray!5, colframe=gray!30, arc=2pt, boxrule=0.6pt,
  left=8pt, right=8pt, top=6pt, bottom=6pt, title={\textbf{Corollary 1}}]
For the dual-axis system (Eq.~\ref{eq:dual_axis}), increasing domain routing accuracy
$\Pr[d^*(q) = d_{\text{gold}}(q)]$ toward its upper bound does not necessarily
increase expected generation quality $\mathbb{E}[Q(\pi_{DF}(q), q)]$.
\end{tcolorbox}

\paragraph{Reasoning.}
Quality depends on three jointly determined components:
(i) adapter training distribution match, (ii) retrieval partition relevance, and
(iii) scaffold alignment. Even when domain routing is oracle-correct, if the
oracle-selected adapter $e_{d_{\text{gold}}}$ has a worse training distribution match
for a given query than a non-gold adapter, or if the oracle partition retrieves
less relevant passages than the predicted partition, quality may decline.
Formally, the quality function $Q$ is not monotone in domain accuracy alone.

\paragraph{Empirical verification.}
Table~\ref{tab:domainrouter} (domain router necessity experiment) shows that
oracle-domain routing degrades ROUGE-L by $-1.72$~pp overall and $-1.96$~pp
single-domain relative to NexRoute predicted routing ($p < 0.001$ for both).
This empirically instantiates Corollary~1 and is the central unexpected finding
of this work: optimizing routing accuracy is not the productive design choice.

\paragraph{Implication for system design.}
Corollaries~1 and Proposition~1 together shift the design objective from
``maximize routing accuracy'' to ``maximize structural alignment between routing
decisions and downstream system resources.'' This reframing motivates NexRoute's
architecture: the nexus point selects not just the expert but the evidence scope
and reasoning format simultaneously, treating structural binding as the primary
design variable.

%%==============================================================================
%% NOTE FOR REVISION: Section 3 (Methodology) should be restructured as:
%% 3.0 Theoretical Motivation (this section)
%% 3.1 Problem Formulation and Framework Overview  (existing Section 3.1)
%% 3.2 Domain-Partitioned Knowledge Base          (existing Section 3.2)
%% ... (rest unchanged)
%%
%% This places the formal analysis BEFORE the architecture description,
%% establishing "why" before "what", which directly addresses the novelty concern.
%%==============================================================================
